\chapter{Go-To-Market Approach}

The motion should be simple and realistic.
We are not running a broad campaign.
We are not waiting for the market to discover us.
We are going directly to the right people with the right message.

\section{Motion}

\begin{itemize}
  \item Founder-led.
  \item Direct outreach.
  \item Network-based.
\end{itemize}

This is the right motion because the wedge is still being validated, the message is still being refined, and founder conversations generate the fastest learning at this stage.

\section{How the motion works}

\subsection{Step 1 --- Identify and reach the right teams}

We target teams that already fit the wedge profile from Chapter~5.

Outreach should be direct and specific.
Not ``we built a powerful AI platform.''

Instead:

\begin{insightbox}{Outreach framing that works}
``Are you using AI to take actions in support or ops?

When it makes the wrong call, can you explain why?

If not, we should talk.''
\end{insightbox}

This framing pre-qualifies the conversation and makes the pain immediately visible to the right buyer.

\subsection{Step 2 --- Sell one workflow, not the full vision}

The first sale is:

\begin{itemize}
  \item one workflow,
  \item one pain point,
  \item one pilot.
\end{itemize}

This keeps conversations manageable, keeps the buyer's internal champion in control, and keeps us focused on proving value rather than educating.

Selling the full platform vision in the first conversation almost always slows things down.

\subsection{Step 3 --- Run the pilot and land the evidence}

After the pilot succeeds:

\begin{itemize}
  \item document what was shown,
  \item get a clear statement from the customer,
  \item use it in the next conversation.
\end{itemize}

One strong pilot outcome is worth more than ten vague conversations.

\subsection{Step 4 --- Expand from pilot to commercial relationship}

After a successful pilot, expansion happens in one of three ways:

\begin{enumerate}
  \item More workflows in the same team.
  \item More teams in the same company.
  \item A reference story that opens doors at similar companies.
\end{enumerate}

This is where Connector becomes sticky.
Once a company relies on us to make AI decisions explainable and defensible, we become part of how they run AI safely.

\section{Why founder-led is correct right now}

At this stage, the founder should be in most early conversations because:

\begin{itemize}
  \item the wedge language is still being sharpened,
  \item objections reveal what needs to change in positioning,
  \item the pilot design benefits from direct customer input,
  \item early buyers often want to know who is building this.
\end{itemize}

Founder-led is not a weakness.
It is the right instrument for the validation phase.

\section{Phase 1 targets for 2026}

By end of Q2 2026:

\begin{itemize}
  \item 3--5 qualified conversations with teams in the target segments,
  \item 1--2 paid pilots in progress or completed,
  \item at least one clear pilot win documented with a customer statement.
\end{itemize}

By end of Q3 2026:

\begin{itemize}
  \item 2--3 pilots completed with documented outcomes,
  \item at least 1 expansion conversation underway,
  \item a tightened message based on what objections surfaced most often.
\end{itemize}

These are not aggressive targets.
They are honest targets for a founder-led wedge validation phase.
